% Options for packages loaded elsewhere
\PassOptionsToPackage{unicode}{hyperref}
\PassOptionsToPackage{hyphens}{url}
%
\documentclass[
  11pt,
]{article}
\usepackage{amsmath,amssymb}
\usepackage{iftex}
\ifPDFTeX
  \usepackage[T1]{fontenc}
  \usepackage[utf8]{inputenc}
  \usepackage{textcomp} % provide euro and other symbols
\else % if luatex or xetex
  \usepackage{unicode-math} % this also loads fontspec
  \defaultfontfeatures{Scale=MatchLowercase}
  \defaultfontfeatures[\rmfamily]{Ligatures=TeX,Scale=1}
\fi
\usepackage{lmodern}
\ifPDFTeX\else
  % xetex/luatex font selection
\fi
% Use upquote if available, for straight quotes in verbatim environments
\IfFileExists{upquote.sty}{\usepackage{upquote}}{}
\IfFileExists{microtype.sty}{% use microtype if available
  \usepackage[]{microtype}
  \UseMicrotypeSet[protrusion]{basicmath} % disable protrusion for tt fonts
}{}
\makeatletter
\@ifundefined{KOMAClassName}{% if non-KOMA class
  \IfFileExists{parskip.sty}{%
    \usepackage{parskip}
  }{% else
    \setlength{\parindent}{0pt}
    \setlength{\parskip}{6pt plus 2pt minus 1pt}}
}{% if KOMA class
  \KOMAoptions{parskip=half}}
\makeatother
\usepackage{xcolor}
\usepackage[margin=1in]{geometry}
\usepackage{graphicx}
\makeatletter
\def\maxwidth{\ifdim\Gin@nat@width>\linewidth\linewidth\else\Gin@nat@width\fi}
\def\maxheight{\ifdim\Gin@nat@height>\textheight\textheight\else\Gin@nat@height\fi}
\makeatother
% Scale images if necessary, so that they will not overflow the page
% margins by default, and it is still possible to overwrite the defaults
% using explicit options in \includegraphics[width, height, ...]{}
\setkeys{Gin}{width=\maxwidth,height=\maxheight,keepaspectratio}
% Set default figure placement to htbp
\makeatletter
\def\fps@figure{htbp}
\makeatother
\setlength{\emergencystretch}{3em} % prevent overfull lines
\providecommand{\tightlist}{%
  \setlength{\itemsep}{0pt}\setlength{\parskip}{0pt}}
\setcounter{secnumdepth}{5}
\usepackage{booktabs}
\usepackage{threeparttable}
\usepackage{amsmath}
\usepackage{graphicx}
\usepackage{float}
\usepackage{caption}
\captionsetup[figure]{labelformat=default,labelsep=period,name=Figure}
\ifLuaTeX
  \usepackage{selnolig}  % disable illegal ligatures
\fi
\IfFileExists{bookmark.sty}{\usepackage{bookmark}}{\usepackage{hyperref}}
\IfFileExists{xurl.sty}{\usepackage{xurl}}{} % add URL line breaks if available
\urlstyle{same}
\hypersetup{
  pdftitle={How Do Users Review Spotify? Analyzing Emotion, Topics, and Platform Replies Using Text-as-Data},
  pdfauthor={Yihan Chen},
  hidelinks,
  pdfcreator={LaTeX via pandoc}}

\title{How Do Users Review Spotify? Analyzing Emotion, Topics, and
Platform Replies Using Text-as-Data}
\author{Yihan Chen}
\date{2025-05-03}

\begin{document}
\maketitle

{
\setcounter{tocdepth}{3}
\tableofcontents
}
\begin{center}
\textbf{Georgetown University, McCourt School of Public Policy}  
Washington, DC 20001  
\end{center}

\section*{Significance Statement}

Understanding how Spotify selectively replies to user reviews sheds
light on platform governance, digital fairness, and customer engagement
strategy. By analyzing emotional tone and review content using TAD and
logistic regression, this study highlights patterns in review topics and
corporate responsiveness, contributing to business insight and broader
discussions on digital transparency and priority.

\section*{Abstract}

This study analyzes 60,000 Spotify reviews to examine what drives
platform replies. Using text classification, topic modeling, and
logistic regression, we find replies are more likely for detailed,
negative feedback---especially with technical issues or emotions like
disgust. Joy, praise, and competitor mentions have little impact.

\begin{center}\rule{0.5\linewidth}{0.5pt}\end{center}

\hypertarget{introduction}{%
\section{Introduction}\label{introduction}}

\hypertarget{motivation}{%
\subsection{Motivation}\label{motivation}}

Spotify has become the global leader in music streaming, with over 252
million premium subscribers as of Q3 2024 (Statista, 2024; Figure 1).

\begin{figure}

{\centering \includegraphics[width=0.8\linewidth]{figures/subscriber} 

}

\caption{Number of Spotify premium subscribers worldwide from Q1 2015 to Q4 2024.}\label{fig:fig.subs}
\end{figure}

While subscription metrics reveal growth, understanding user
satisfaction and retention requires deeper analysis beyond numerical
figures. User-generated content (UGC), especially app store reviews,
provides insights into customer experiences.

These reviews reflect not only product satisfaction but also user
frustrations and unmet expectations. Platforms like Spotify occasionally
respond to user reviews, but the dynamics of which reviews receive
replies remain largely unexplored.

\hypertarget{prior-research}{%
\subsection{Prior Research}\label{prior-research}}

Literature supports that text analysis informs consumer insights. Pang
\& Lee (2008) and Liu (2012) laid the foundation of sentiment analysis
for understanding customer satisfaction. Tirunillai \& Tellis (2012)
showed that user-generated content predicts business performance, while
Blei et al.~(2003) introduced LDA topic modeling to identify consumer
concerns.

In digital media, Datta et al.~(2018) found that users' sentiments
correlate with music subscriber growth, and Piskorski et al.~(2019)
demonstrated that negative sentiment spikes signal churn risk. These
findings support the analysis of user sentiment in Spotify's reviews and
can inform market decisions.

\hypertarget{research-questions}{%
\subsection{Research Questions}\label{research-questions}}

\begin{itemize}
\tightlist
\item
  \textbf{RQ1}: How do Spotify users describe their experience in
  written reviews?\\
\item
  \textbf{RQ2}: Which review topics or tones are more likely to receive
  an official reply---and what does this reveal about platform
  priorities or transparency?
\end{itemize}

\hypertarget{report-overview}{%
\subsection{Report Overview}\label{report-overview}}

This report begins by describing the dataset and annotation process,
followed by supervised classification and topic modeling. I then
evaluate emotional tone and reply likelihood through logistic
regression, and conclude with a discussion of key findings and
implications.

\hypertarget{data-and-methods}{%
\section{Data and Methods}\label{data-and-methods}}

\hypertarget{data-source-and-unit-of-observation}{%
\subsection{Data Source and Unit of
Observation}\label{data-source-and-unit-of-observation}}

The dataset, published on Kaggle (2022), contains 61,594 Spotify app
reviews from the Google Play Store. Each entry includes the text of the
review, timestamp, rating (1--5), helpfulness score, and reply (if any).
Each observation corresponds to one user's review.

\hypertarget{variables-of-interest}{%
\subsection{Variables of Interest}\label{variables-of-interest}}

The dataset used in this study comprises a set of variables that support
both classification and regression analysis (Table 1). Each review is
accompanied by structured metadata and derived features that inform our
analysis of user experience and reply behavior.

\begin{center}\includegraphics[width=0.8\linewidth]{figures/table1} \end{center}

\hypertarget{data-wrangling-steps}{%
\subsection{Data Wrangling Steps}\label{data-wrangling-steps}}

To prepare the dataset for analysis, I implemented a structured data
cleaning and transformation pipeline.

First, I converted all review text to lowercase, removed duplicates, and
eliminated special characters to ensure consistency. The cleaned text
was then tokenized and matched with the NRC lexicon to compute
review-level emotion scores across ten emotion categories.

Using a manually labeled random sample of 2,000 reviews, I initially
created four functional categories: Technical Issues, UX \& Features,
Price \& Ads, and Other. However, due to overlapping themes and
classification ambiguity---especially between Technical and UX
feedback---I consolidated these into a binary scheme to improve
annotation consistency and model performance. Therefore, the final
categories were: (1) ``App Experience'': reviews mainly focused on
technical problems, user interface, and app functionality; (2)
``Other'': reviews related to pricing, ads, music content, or general
sentiment. The ``App Experience'' category may include all forms of
interaction with the Spotify platform, including the mobile app,
website, and other interfaces. It may also cover reviews that partially
mention pricing, praise, or subscription but are primarily focused on
user experience. This simplified structure allowed the classifier to
more reliably distinguish actionable user feedback tied to user
experience.

I also engineered several features to support downstream analysis.
Binary flags were created to indicate whether a review received an
official Spotify reply and whether a competing platform (e.g., Apple
Music) was mentioned. I computed word counts for each review and
detected specific competitor references.

For classification, I constructed both TF-IDF and word embedding
matrices to predict topic labels. To perform unsupervised topic
modeling, I filtered the App Experience reviews and preprocessed them
for input into an LDA model.

To address class imbalance in reply behavior, I applied SMOTE (Synthetic
Minority Oversampling Technique), generating synthetic samples for the
minority class (Replied = 1). This rebalancing allowed me to fit more
robust logistic regression models.

\hypertarget{analysis}{%
\section{Analysis}\label{analysis}}

\hypertarget{supervised-classification}{%
\subsection{Supervised Classification}\label{supervised-classification}}

To distinguish between App Experience and Other reviews, I trained three
supervised models---Na"\{i\}ve Bayes, SVM, and Lasso---on the annotated
subset. Among these, the Lasso logistic regression model achieved the
highest accuracy (83\%) and interpretability (Table 2). I applied this
model to classify the remaining reviews in the dataset.

\begin{center}\includegraphics[width=0.8\linewidth]{figures/table2} \end{center}

\hypertarget{topic-modeling}{%
\subsection{Topic Modeling}\label{topic-modeling}}

LDA was applied to App Experience reviews (\textasciitilde38,000
records), revealing six subtopics (Table 3): Technical Fixes, Account \&
Subscription, General UX Feedback, Playback Issues, Music \& Playlist
Features, and Positive Experience.

\begin{center}\includegraphics[width=0.8\linewidth]{figures/table3} \end{center}

\hypertarget{sentiment-and-emotion}{%
\subsection{Sentiment and Emotion}\label{sentiment-and-emotion}}

Sentiment was categorized based on ratings: 1--2 as negative, 3 neutral,
and 4--5 positive. NRC Lexicon provided 10 emotion scores (e.g., anger,
joy, disgust) for each review.

\hypertarget{logistic-regression-reply-prediction}{%
\subsection{Logistic Regression: Reply
Prediction}\label{logistic-regression-reply-prediction}}

I modeled reply likelihood using emotional scores, review length,
competitor mentions, and topic labels. A second model applied SMOTE to
account for class imbalance. The model specification is shown as
follows:

\begin{center}\includegraphics[width=0.8\linewidth]{figures/logit} \end{center}

\hypertarget{results}{%
\section{Results}\label{results}}

\hypertarget{descriptive-overview}{%
\subsection{Descriptive Overview}\label{descriptive-overview}}

49 percent of Spotify reviews are positive in tone, with around 60
percent classified under the App Experience category. However, only 3.5
percent of reviews received an official reply from the platform,
highlighting a significant gap in engagement (Figure 2).

\begin{figure}

{\centering \includegraphics[width=0.8\linewidth]{figures/plot123} 

}

\caption{Descriptive overview: sentiment tone and App Experience classification share among Spotify reviews.}\label{fig:fig.2}
\end{figure}

\hypertarget{sentiment-emotion-score-and-reply}{%
\subsection{Sentiment, Emotion Score, and
Reply}\label{sentiment-emotion-score-and-reply}}

Spotify tends not to respond to emotionally extreme or overtly positive
reviews; instead, it is more likely to engage with detailed expressions
of dissatisfaction. Reviews expressing disgust or positive sentiment are
more likely to receive a response (Figure 3), with disgust showing a
significant positive effect (coefficient = 0.464, \(p < 0.001\)). In
contrast, reviews characterized by anger, joy, or trust tend to be
overlooked---joy (coefficient = -0.240, \(p < 0.05\)) and trust
(coefficient = -0.226, \(p < 0.05\)) both significantly reduce the
likelihood of a reply.

\begin{figure}

{\centering \includegraphics[width=0.8\linewidth]{figures/emotionsreply} 

}

\caption{Average emotion scores by reply status: replied vs. unreplied reviews across ten NRC emotion categories.}\label{fig:fig.3}
\end{figure}

\hypertarget{topics-spotify-prioritizes-problem-solving}{%
\subsection{Topics: Spotify Prioritizes
Problem-Solving}\label{topics-spotify-prioritizes-problem-solving}}

Spotify appears more responsive to critical reviews tied to app
performance than to general praise. Approximately 51\% of reviews
classified as App Experience were negative, compared to just 23\% in the
Other category. Furthermore, over 75\% of all replies were directed to
App Experience reviews.

Different topics share broadly similar distributions of emotion scores
(Figure 4), with one key exception: the Positive Experience topic tends
to reflect high levels of joy and positive sentiment, yet is one of the
least likely to receive a reply.

Spotify prioritizes resolving technical access and service issues over
addressing feature feedback or acknowledging general satisfaction.
Compared to the reference category (Account \& Subscription), reviews
classified as Music and Playlist Features (coefficient = -0.717,
\(p < 0.05\)), Other (coefficient = -1.031, \(p < 0.001\)), and Positive
Experience (coefficient = -0.915, \(p < 0.01\)) are significantly less
likely to receive a response.

\begin{figure}

{\centering \includegraphics[width=0.8\linewidth]{figures/emotionstopic} 

}

\caption{Emotion scores by topic: each LDA-derived topic displays its corresponding average emotion levels using the NRC lexicon.}\label{fig:fig.4}
\end{figure}

\hypertarget{competitor-mentions-mentioning-others-doesnt-get-noticed}{%
\subsection{Competitor Mentions: Mentioning Others Doesn't Get
Noticed}\label{competitor-mentions-mentioning-others-doesnt-get-noticed}}

The ten most frequently mentioned competitors in Spotify reviews are
Pandora, YouTube Music, Apple Music, Amazon Music, SoundCloud, Tidal,
Deezer, Gaana, Wynk, and Napster. Among the top 5 mentioned competitors,
Apple Music, YouTube Music, and Amazon Music tend to show the most joy,
positive sentiment, and sadness scores (Figure 5). SoundCloud, YouTube
Music, and Amazon Music are often mentioned in negative contexts,
whereas Apple Music and Pandora are more frequently cited in positive or
neutral reviews).

\begin{figure}

{\centering \includegraphics[width=0.8\linewidth]{figures/competitor} 

}

\caption{Emotion scores by top 5 competitors}\label{fig:fig.5}
\end{figure}

However, Spotify did not explicitly engage with replies containing
competitors. According to the logistic regression model, the presence of
a competitor mention is not a statistically significant predictor of
receiving a reply (coefficient = -0.103, \(p > 0.05\)).

\hypertarget{word-count-spotify-values-detailed-user-feedback}{%
\subsection{Word Count: Spotify Values Detailed User
Feedback}\label{word-count-spotify-values-detailed-user-feedback}}

Most reviews are relatively concise: the average review contains 32
words, and the interquartile range spans from 14 to 43 words, though
some extend up to 738. Reviews that are longer in length are
significantly more likely to receive a response: each additional word
increases the probability of receiving a reply by approximately 1\%
(coefficient = 0.010, \(p < 0.01\)).

\hypertarget{smote-adjustment-for-class-imbalance}{%
\subsection{SMOTE Adjustment for Class
Imbalance}\label{smote-adjustment-for-class-imbalance}}

Spotify's engagement is systematically focused on complaints, even under
balanced sampling. To validate the robustness of these findings, I
re-estimated the logistic regression using a SMOTE-balanced dataset to
address the low reply rate. The results were largely consistent:
emotionally negative tones like disgust and technical topics like
Account \& Subscription remained strong predictors of reply likelihood
(Figure 6; Table 4).

\begin{figure}

{\centering \includegraphics[width=0.8\linewidth]{figures/comparison} 

}

\caption{Comparison of Coefficient Estimates}\label{fig:fig.6}
\end{figure}

\begin{center}\includegraphics[width=0.8\linewidth]{figures/table4} \end{center}

\hypertarget{discussion}{%
\section{Discussion}\label{discussion}}

Spotify replies to only a small fraction of user reviews and does so
selectively, prioritizing longer, more detailed reviews that express
dissatisfaction, especially those involving technical or service-related
issues. Reviews with emotionally nuanced tones like disgust are more
likely to receive a response, while those expressing joy, trust, or
general praise tend to be overlooked. Topic modeling and regression
results further confirm that Spotify is most responsive to reviews under
categories like Account \& Subscription. Mentioning competitors such as
Apple Music or YouTube Music does not significantly influence reply
behavior.

From a business strategy perspective, the analysis identifies key
experience themes and sentiment patterns that can inform improvements in
Spotify's app functionality, user satisfaction, and retention efforts.
By highlighting which types of reviews prompt engagement---namely those
that are detailed, emotionally nuanced, and related to technical
service---Spotify can better prioritize customer service efforts and
refine its product experience.

From a governance perspective, this study provides transparency into how
Spotify manages public engagement through app store replies. It reveals
selective responsiveness to certain emotional tones and topics, raising
broader questions about fairness, accountability, and communication
standards in platform behavior. The use of text-as-data methods
contributes to ongoing policy discussions by empirically examining how
platform design choices affect user interaction and perceived
responsiveness.

Future research could expand this analysis by incorporating data from
other platforms (e.g., Apple or YouTube Music) to allow for comparative
studies. Further improvements could also include training deep learning
models on a larger, more diverse dataset or integrating user metadata to
explore demographic patterns in reply behavior. Besides, while reply
patterns suggest engagement priorities, further work is needed to fully
assess transparency and policy implications.

\hypertarget{limitations}{%
\subsection{Limitations}\label{limitations}}

While the findings offer clear patterns, several limitations remain.
First, TF-IDF outperformed word embeddings in classification, which may
reflect insufficient embedding tuning or the brevity of most reviews.

Second, topic modeling and classification outputs showed some conceptual
overlap, limiting interpretability. Enhanced preprocessing or
alternative models could improve topic separation.

Third, although the number of retained duplicate reviews is small (270),
their short and repetitive nature may reduce analytical quality. Word
clouds also revealed minor stemming issues, such as splitting ``song''
and ``songs.''

Finally, fewer than 3.5\% of reviews received a reply, significantly
limiting the statistical power of the regression model.

\hypertarget{r-packages-used}{%
\subsection{R Packages Used}\label{r-packages-used}}

\begin{itemize}
\tightlist
\item
  \textbf{Data Cleaning and Transformation}: \texttt{tidyverse},
  \texttt{janitor}, \texttt{textclean}, \texttt{hunspell}\\
\item
  \textbf{Text Preprocessing}: \texttt{tidytext}, \texttt{quanteda},
  \texttt{textmineR}, \texttt{quanteda.textstats}\\
\item
  \textbf{Topic Modeling}: \texttt{topicmodels}, \texttt{broom.mixed}\\
\item
  \textbf{Classification}: \texttt{caret}, \texttt{glmnet},
  \texttt{e1071}, \texttt{naivebayes}\\
\item
  \textbf{Visualization}: \texttt{ggplot2}, \texttt{patchwork},
  \texttt{wordcloud}, \texttt{RColorBrewer}\\
\item
  \textbf{Model Reporting}: \texttt{pscl}, \texttt{modelsummary},
  \texttt{gridExtra}\\
\item
  \textbf{SMOTE Handling}: \texttt{smotefamily}
\end{itemize}

\hypertarget{references}{%
\section{References}\label{references}}

\begin{itemize}
\item
  Looka. (2024). \emph{The Spotify Logo: Embodying the sonic style of a
  streaming pioneer}. Retrieved April 2025, from
  \url{https://looka.com/blog/spotify-logo/}
\item
  Statista. (2024). \emph{Number of Spotify premium subscribers
  worldwide from 1st quarter 2015 to 4th quarter of 2024}. Retrieved
  from
  \url{https://www.statista.com/statistics/244995/number-of-paying-spotify-subscribers/}
\item
  Pang, B., \& Lee, L. (2008). \emph{Opinion mining and sentiment
  analysis}. Foundations and Trends® in Information Retrieval, 2(1--2),
  1--135. \url{https://www.nowpublishers.com/article/Details/INR-011}
\item
  Liu, B. (2012). \emph{Sentiment analysis and opinion mining}.
  Springer.
  \url{https://link.springer.com/book/10.1007/978-3-031-02145-9}
\item
  Tirunillai, S., \& Tellis, G. J. (2012). Does online chatter really
  matter? Dynamics of user-generated content and stock performance.
  \emph{Marketing Science}, 31(2), 198--215.
  \href{https://www.researchgate.net/publication/228151131_Does_Online_Chatter_Really_Matter_Dynamics_of_User-Generated_Content_and_Stock_Performance}{https://www.researchgate.net/publication/228151131}
\item
  Blei, D. M., Ng, A. Y., \& Jordan, M. I. (2003). Latent Dirichlet
  Allocation. \emph{Journal of Machine Learning Research}, 3, 993--1022.
  \href{https://www.researchgate.net/publication/221620547_Latent_Dirichlet_Allocation}{https://www.researchgate.net/publication/221620547}
\item
  Datta, H., Knox, G., \& Bronnenberg, B. J. (2018). Changing their
  tune: How consumers' adoption of online streaming affects music
  consumption and discovery. \emph{Marketing Science}, 37(1), 5--21.
  \href{https://www.researchgate.net/publication/293753759_Changing_Their_Tune_How_Consumers\%27_Adoption_of_Online_Streaming_Affects_Music_Consumption_and_Discovery}{https://www.researchgate.net/publication/293753759}
\item
  Piskorski, M. J., Shin, J., \& Smith, H. (2019). Social media and
  customer retention: Implications for the luxury beauty industry.
  \emph{Journal of Business Research}, 95, 211--219.
  \href{https://www.researchgate.net/publication/275947382_Social_Media_and_Customer_Retention_Implications_for_the_Luxury_Beauty_Industry}{https://www.researchgate.net/publication/275947382}
\end{itemize}

\end{document}
